\section{Step 46 Count lengths without treating zero as a code}\label{png-huff-s46}
Fill \texttt{huff\_tree\_count}. Read each length as an unsigned byte and reject values above fifteen before indexing the count array. Increment its corresponding word count. Separately count nonzero lengths and retain their maximum. Zero belongs in the histogram but does not increase the number of present codes.

For the first fixed table, the nonzero counts are 24 at length seven, 152 at eight and 112 at nine. There are 288 present symbols and maximum length nine. A later distance alphabet may contain only absent entries; its policy decision needs the present count, not the total number of declared slots.

\paragraph{Observe} Use \texttt{s46-1}; stop before \texttt{huff\_tree\_counted}. Inspect the state described above before continuing.

\Needspace{7\baselineskip}
\paragraph{Check} Does a zero length denote a zero-bit data code?
\begin{itemize}
  \item[A.] Yes, one that produces a zero byte.
  \item[B.] Only in the fixed alphabet.
  \item[C.] No; that symbol is absent.
\end{itemize}
